\section{讨论}

\subsection{现有工具的缺陷不在税率，而在工具形式}

把第~\ref{sec:magnitude}、\ref{sec:identification}、\ref{sec:comparison}~节放在一起，得到一个明确的判断：\textbf{消费税在这个问题上的问题不是“太低”，而是“不对准”。}

一个全国统一、按年适用、在生产环节征收的从价税，其全部信息含量是“全国平均而言，一件烟花应该贵 15\%”。而本文测出的损害的信息含量是“某个城市在某一年，浓度上升了几十到几百微克”。第~\ref{sec:variance}~节的方差分解把这句话变成了数字：城市间变异占 22.0\%，年份间占 12.7\%，\textbf{65.3\% 落在城市 $\times$ 年份这一维度上}。\textbf{两者之间不是量的差距，是维度的差距。}

这也解释了为什么“提高税率”的主张在实践上很难落地：把税率从 15\% 提到 30\%，全国所有城市、全年 365 天一起变贵，而损害只集中在少数城市的十几个小时。\textbf{代价是普遍的，收益是局部的}——这正是 Oates \& Baumol（1975）所说的可见性不对称：设定配额时可以严格而不必面对成本，而提高税率则必须面对每一个纳税人。

\subsection{如果要改，改什么}

按本文的逻辑，可能的方向有三个，按可行性递减：

\begin{enumerate}
  \item \textbf{从价改从量。} 按净火药量定额征收，使税负与排放成比例。这是唯一不需要改变工具层级的改进，也最容易被现行征管吸收。
  \item \textbf{在年度之内做时段差异化。} 例如对春节前若干日的销售适用更高税率。这会把时间维度从“零分辨率”提高到“粗分辨率”，但仍然无法区分城市。
  \item \textbf{把工具层级下放。} 允许地方在国家标准之上设定附加税率。这会直接冲击消费税作为中央税的性质，在现行体制下几无可能。
\end{enumerate}

本文不主张其中任何一个。列出来的目的是说明：\textbf{在这三个方向上，“提高税率”甚至不在其中。}

\subsection{本文的边界}\label{sec:limits}

以下六点必须与结论同时出现，否则会被合理地批评为过度解读：

\begin{enumerate}
  \item \textbf{第一层的因果识别失败，不等于禁放无效。} 本文的失败是“无法分离”，不是“效应为零”。
  \item \textbf{庇古校准不报单点。} 比值跨约三个数量级，只报区间与中位数，且中位数必须与区间同时出现。
  \item \textbf{“未找到”不等于“没有”。} 34 城中有 4 城文件未找到，正文以取得原文的 30 城为分母。
  \item \textbf{牵头单位各地不一}，不得写成“由公安局牵头”。
  \item \textbf{春节是短期暴露}，不得套用全年长期剂量反应系数。
  \item \textbf{罚款与税收是转移支付}，不计入社会净成本。
\end{enumerate}

此外还有两处方法上的自我限制。其一，本文的机制证据（\SOtwo{} 同向、\NOtwo{} 反向）依赖污染物结构，但没有做气象匹配（边界层高度、风速）；在静稳天气下，任何一次排放增加都会产生更大的浓度响应，ERA5 气象匹配是本项目已识别但尚未完成的加固工作。其二，第~\ref{sec:variance}~节的方差分解使用已实现的尖峰增量，其中含不可预测的气象成分，因此 65.3\% 只应读作分辨率缺口的量级。

\subsection{一项方法论教训}

本文第一节曾写出一个 $t=-4.5$ 的结果。它通过了常规的显著性检验，通过了随机化推断的自校准，也在两个设定中的一个成立。它是假的。

发现它的是三件事：\textbf{安慰剂“日”扫描}（效应出现在不该出现的地方）、\textbf{有效聚类数}（306 个聚类里只有 17 个真正贡献方差）、以及\textbf{量级先验}（效应大于中位城市的整个尖峰）。

\textbf{这三件事都不是显著性检验。} 它们的共同点是：不问“这个系数有多大、多显著”，而问“这个系数在物理上可能吗、它从哪里来”。这是本文愿意留下的唯一一条方法建议。
